\LectureTitle{그림의 규칙을 양자회로로 옮기는 법}{아다마르·파울리 기저에서 관측량과 측정까지}{한국어 전체 개정판 v3 · 2026년 11월 28일\\책임 편집자 김현중 · 편집부(AEDAUT) · OpenAI Codex 보조}

\section{네 칸 그림에서 출발하기}

흰 칸 세 개와 검은 칸 하나가 있는 작은 그림을 생각해 보자. 칸의 주소를
\(00,01,10,11\)로 정하고, 흰색은 0, 검은색은 1로 적으면 그림은
\[
f=(0,0,0,1)
\]
이라는 표가 된다. 두 큐비트에는 이 네 주소에 대응하는 상태
\(\lvert00\rangle,\lvert01\rangle,\lvert10\rangle,\lvert11\rangle\)가 있다.
\(\lvert x\rangle\)는 주소 \(x\)의 기본 상태이고, \(f(x)\)는 그 칸의 0 또는 1이다.
여러 상태를 더해 쓸 때 각 상태 앞에 곱하는 수를 진폭이라고 부른다.
이제 검은 칸에 해당하는 진폭에만 마이너스 부호를 붙이는 규칙을 만들어 보자.
\[
U_f\lvert x\rangle=(-1)^{f(x)}\lvert x\rangle,
\qquad U_f=\operatorname{diag}(1,1,1,-1).
\]
\(f(00)=0\)이면 \((-1)^0=+1\), \(f(11)=1\)이면 \((-1)^1=-1\)을 곱한다.
\(\operatorname{diag}\)는 괄호 안의 수를 행렬의 대각선에 놓으라는 뜻이다.
이 규칙은 그림 자체를 양자 상태로 준비하는 일이 아니다. 어떤 입력 상태가 들어오든,
주소 \(11\)에 붙은 진폭의 부호를 바꾸는 \textbf{위상 오라클}을 정한 것이다.

이 노트에서는 이 작은 규칙을 실제 게이트 회로로 옮기는 과정을 따라간다.
그 과정에 필요한 파울리 기저를 이해하려면, 먼저 상태를 적는 벡터와 상태에 작용하는 행렬을
구별해야 한다. 여기에 IBM의 관측량 문서~\cite{ibm-observables}에서 다루는
측정과 기대값을 연결했다. 행렬을 파울리 항으로 나누면 회로를 만드는 데에도,
측정값을 계산하는 데에도 쓸 수 있다. 두 쓰임에서 무엇을 더하고 무엇을 실행하는지
작은 예제로 하나씩 확인하자.

\begin{notebox}{읽는 데 필요한 약속}
\(\lvert\psi\rangle\)는 상태를 나타내는 열벡터이고, \(i\)는 \(i^2=-1\)인 수다.
\(\dagger\)는 행과 열을 바꾸면서 복소수를 켤레로 바꾸는 연산이다.
켤레는 \(a+bi\)의 \(i\) 앞 부호를 뒤집어 \(a-bi\)로 바꾸는 것이다(\(a,b\)는 실수).
예를 들어 \((a,b)^\mathsf{T}\)의 \(\dagger\)는 \((a^*,b^*)\)다.
두 큐비트의 주소는 끝까지 \(\lvert q_1q_0\rangle\) 순서로 쓴다.
벡터 성분과 행렬의 행·열 순서는 \(00,01,10,11\)이다.
\end{notebox}

\section{상태에는 진폭이 있고, 측정에는 결과가 있다}

한 큐비트의 가장 기본적인 두 상태는
\[
\lvert0\rangle=\begin{pmatrix}1\\0\end{pmatrix},\qquad
\lvert1\rangle=\begin{pmatrix}0\\1\end{pmatrix}
\]
이다. 일반적인 순수 상태는
\(\lvert\psi\rangle=a\lvert0\rangle+b\lvert1\rangle\)로 적는다.
\(a,b\)를 진폭이라고 부른다. 여기서 기본 측정은 0과 1을 구별하는 계산 기저 측정이다.
이 측정에서 0과 1이 나올 확률은 각각
\(\lvert a\rvert^2,\lvert b\rvert^2\)이고, 두 확률의 합이 1이 되도록
\(\lvert a\rvert^2+\lvert b\rvert^2=1\)을 요구한다.

다음 두 상태를 나란히 놓으면 진폭과 확률의 차이가 드러난다.
\[
\lvert+\rangle=\frac{\lvert0\rangle+\lvert1\rangle}{\sqrt2},\qquad
\lvert-\rangle=\frac{\lvert0\rangle-\lvert1\rangle}{\sqrt2}.
\]
어느 상태를 기본 측정으로 읽어도 0과 1의 확률은 반반이다. 그러나 두 상태는 서로 다르다.
두 번째 진폭의 상대적인 부호가 다르기 때문이다. 이 차이는 알맞은 게이트를 거친 뒤
측정하면 드러난다. 따라서 위상 오라클이 부호를 바꿨다는 사실만으로 기본 측정의
확률이 곧바로 달라졌다고 말할 수는 없다.

\subsection{관측량은 어떤 값을 읽을지를 정하는 행렬이다}

기본 측정의 비트에 0이면 \(+1\), 1이면 \(-1\)을 붙여 읽어 보자.
이 측정에 대응하는 관측량은
\[
Z=\begin{pmatrix}1&0\\0&-1\end{pmatrix}
\]
다. \(Z\lvert0\rangle=+\lvert0\rangle\),
\(Z\lvert1\rangle=-\lvert1\rangle\)이므로 두 상태는 측정값이 확실히 정해진 상태다.
이처럼 행렬을 적용했을 때 같은 벡터에 수만 곱해지는 벡터를 \textbf{고유벡터},
그 수를 \textbf{고윳값}이라고 부른다.

관측량 \(O\)는 \(O^\dagger=O\)인 에르미트 행렬로 나타낸다.
이 성질은 측정 가능한 값인 고윳값이 실수가 되게 한다.
한 번의 측정에서는 그 고윳값 중 하나를 얻는다. 같은 상태를 매번 새로 준비해
반복 측정했을 때의 평균을 \textbf{기대값}이라고 부르며,
순수 상태에서는 다음과 같이 계산한다~\cite{preskill}.
\[
\langle O\rangle_\psi=\langle\psi\rvert O\lvert\psi\rangle.
\]
\(Z\)의 경우에는 \(\langle Z\rangle=\lvert a\rvert^2-\lvert b\rvert^2\)다.
\(\lvert0\rangle\)의 평균은 1, \(\lvert+\rangle\)의 평균은 0이다.
평균 0이라는 말은 한 번 측정해서 0이 나왔다는 뜻이 아니다.
이 측정에서 한 번에 나오는 값은 \(+1\) 또는 \(-1\)이다.

\section{파울리 행렬 네 개와 아다마르 게이트}

한 큐비트의 행렬을 표현할 때 특히 편리한 네 행렬이 있다.
\[
I=\begin{pmatrix}1&0\\0&1\end{pmatrix},\quad
X=\begin{pmatrix}0&1\\1&0\end{pmatrix},\quad
Y=\begin{pmatrix}0&-i\\i&0\end{pmatrix},\quad
Z=\begin{pmatrix}1&0\\0&-1\end{pmatrix}.
\]
\(I\)는 상태를 그대로 두고, \(X\)는 \(\lvert0\rangle\)와
\(\lvert1\rangle\)을 맞바꾼다. \(Z\)는 \(\lvert1\rangle\)의 부호만 바꾼다.
\(Y\)는 \(\lvert0\rangle\mapsto i\lvert1\rangle\),
\(\lvert1\rangle\mapsto-i\lvert0\rangle\)로 작용한다.
이 작용을 기억하면 행렬의 부호와 \(i\)가 어디에 있는지 다시 확인할 수 있다.

\lecturefigure{pauli-action}{한 큐비트의 두 기본 상태에 대한 파울리 행렬의 작용. 화살표 옆의 수는 확률이 아니라 진폭에 곱하는 수다. 본문의 네 행렬에서 직접 구성한 설명도다.}

네 행렬을 관측량으로 읽을 수도 있다. \(I\)를 측정하면 언제나 \(+1\)이고,
\(X,Y,Z\)의 가능한 값은 각각 \(\pm1\)이다. 여기서 값은 무차원이다.
실제 스핀 각운동량 \(S_z=(\hbar/2)Z\)를 다룰 때에는 단위와 \(\hbar/2\)가 따로 들어간다.

\subsection{같은 상태에도 서로 다른 측정 질문을 할 수 있다}

\(\lvert+i\rangle=(\lvert0\rangle+i\lvert1\rangle)/\sqrt2\)를 하나 더 준비하자.
이 상태는 \(Y\)의 \(+1\) 고유상태다. 세 상태에서 기대값을 계산하면 다음과 같다.
\begin{center}
\begin{tabular}{lccc}
\toprule
준비한 상태 & \(\langle X\rangle\) & \(\langle Y\rangle\) & \(\langle Z\rangle\)\\
\midrule
\(\lvert0\rangle\) & 0 & 0 & 1\\
\(\lvert+\rangle\) & 1 & 0 & 0\\
\(\lvert+i\rangle\) & 0 & 1 & 0\\
\bottomrule
\end{tabular}
\end{center}
기본 측정에서는 마지막 두 상태가 모두 반반으로 보인다.
\(X\)와 \(Y\) 측정까지 보면 어느 상대위상을 가진 상태인지 구별할 수 있다.
상태를 고정해 두고 어떤 관측량을 읽는지 바꾼 결과다.

\subsection{아다마르 게이트가 바꾸는 것}

아다마르 게이트는
\[
H=\frac1{\sqrt2}\begin{pmatrix}1&1\\1&-1\end{pmatrix},
\qquad H\lvert0\rangle=\lvert+\rangle,\quad
H\lvert1\rangle=\lvert-\rangle
\]
이다. \(H^2=I\)이므로 반대로 \(H\lvert+\rangle=\lvert0\rangle\),
\(H\lvert-\rangle=\lvert1\rangle\)도 성립한다.
앞 절에서 기본 측정으로는 구별하지 못했던 두 부호를, \(H\)를 거친 뒤에는
0과 1로 확실히 구별할 수 있게 된 것이다.

\(H\)의 두 열은 \(X\)의 고유상태 \(\lvert+\rangle,\lvert-\rangle\)다.
그래서 \(H\)는 계산 기저와 \(X\) 고유기저를 연결한다.
이 두 열을 \(H\) 자체의 고유상태라고 부르거나, 모든 관측량의 공통 고유기저라고
부르면 안 된다. 어느 행렬의 고유상태인지를 함께 적어야 뜻이 정확해진다.

\begin{notebox}{게이트와 관측량의 조건은 서로 다르다}
게이트 \(U\)는 \(U^\dagger U=I\)인 유니터리 행렬이다. 상태의 전체 확률을 보존한다.
관측량 \(O\)는 \(O^\dagger=O\)인 행렬이다. 측정할 값과 그 고유상태를 정한다.
\(X,Y,Z,H\)는 두 조건을 모두 만족하지만, 일반적인 관측량이 그대로 게이트가 되는 것은 아니다.
\end{notebox}

\section{벡터의 기저 두 개, 행렬의 기저 네 개}

한 큐비트 상태는 두 성분이므로 \(\lvert0\rangle,\lvert1\rangle\) 두 벡터로 적을 수 있다.
반면 상태에 작용하는 행렬은 \(2\times2\) 크기다. 그 네 자리를 표현할 때
\(I,X,Y,Z\) 네 행렬이 기저가 된다.
\textbf{상태를 펼쳐 쓰는 기준}과 \textbf{연산자를 펼쳐 쓰는 기준}에서 세는 대상이 다른 것이다.

에르미트 행렬 하나를 다음처럼 적어 보자.
\[
O=\begin{pmatrix}u&w-it\\w+it&v\end{pmatrix},
\qquad u,v,w,t\in\mathbb R.
\]
대각 성분 두 개와, 서로 켤레인 비대각 성분의 실수부·허수부가 필요하다.
모두 네 개의 실수다. 네 파울리 행렬을 이용하면 같은 행렬을
\[
O=aI+bX+cY+dZ
 =\begin{pmatrix}a+d&b-ic\\b+ic&a-d\end{pmatrix}
\]
로 적는다. 따라서 \(a=(u+v)/2\), \(d=(u-v)/2\), \(b=w\), \(c=t\)다.
특별한 관측량에만 적용되는 요령이 아니라 모든 한 큐비트 에르미트 행렬을 적는 방법이다.
일반 복소 행렬도 같은 네 행렬로 표현할 수 있으며 그때에는 계수가 복소수일 수 있다.

\lecturefigure{operator-coordinates}{상태를 표현하는 두 벡터와 연산자를 표현하는 네 행렬을 나란히 놓았다. 연산자 쪽의 계수는 행렬의 좌표이며 측정 결과의 확률이 아니다. 독립적인 개수와 본문의 전개식을 바탕으로 구성한 설명도다.}

\begin{examplebox}{직접 만든 관측량 하나를 끝까지 계산하기}
이후 예제에서는
\[
O=2I+\tfrac12X-\tfrac14Y+\tfrac34Z
 =\begin{pmatrix}2.75&0.5+0.25i\\0.5-0.25i&1.25\end{pmatrix}
\]
를 사용한다. \(Y\)의 계수가 음수라서 오른쪽 위 성분의 허수부는 양수가 된다.
앞의 기대값 표를 대입하면
\[
\langle O\rangle_{\lvert0\rangle}=2.75,\qquad
\langle O\rangle_{\lvert+\rangle}=2.5,\qquad
\langle O\rangle_{\lvert+i\rangle}=1.75.
\]
관측량은 같고 준비한 상태가 달라져 평균이 달라졌다.
계수 2, \(1/2\), \(-1/4\), \(3/4\)는 각 측정 평균을 얼마나 반영할지 정하는 가중치다.
확률이 아니므로 음수이거나 1보다 커도 된다.
\end{examplebox}

\section{행렬에서 파울리 계수를 꺼내는 방법}

큐비트가 늘어나면 행렬 성분을 눈으로 비교하기 어려워진다.
이때에는 기저의 직교성을 이용한다. 벡터에서 두 방향의 겹침을 내적으로 계산하듯,
행렬에서는 \(\operatorname{Tr}(A^\dagger B)\)로 겹침을 계산할 수 있다.
이를 힐베르트--슈미트 내적이라고 부른다.
\(\operatorname{Tr}\), 즉 트레이스는 행렬의 대각 성분을 모두 더하는 연산이다.

예를 들어 \(\operatorname{Tr}(I)=2\), \(\operatorname{Tr}(Z)=1-1=0\)이다.
같은 파울리 행렬끼리 곱하면 \(I\)가 되어 트레이스가 2이고,
서로 다른 두 행렬을 곱한 결과의 트레이스는 0이다.
따라서 \(O=aI+bX+cY+dZ\)에 \(X\)를 곱하고 트레이스를 취하면
\[
\operatorname{Tr}(XO)
=a\operatorname{Tr}(X)+b\operatorname{Tr}(X^2)
+c\operatorname{Tr}(XY)+d\operatorname{Tr}(XZ)=2b.
\]
다른 방향의 성분은 사라지고 \(X\) 성분만 남는다. 결국
\[
a=\tfrac12\operatorname{Tr}(O),\quad
b=\tfrac12\operatorname{Tr}(XO),\quad
c=\tfrac12\operatorname{Tr}(YO),\quad
d=\tfrac12\operatorname{Tr}(ZO).
\]
앞의 예제에서는 \(\operatorname{Tr}(O)=4\),
\(\operatorname{Tr}(XO)=1\), \(\operatorname{Tr}(YO)=-1/2\),
\(\operatorname{Tr}(ZO)=3/2\)가 되어 원래 계수를 되찾는다.

\subsection{여러 큐비트에서는 왜 \texorpdfstring{\(1/2^n\)}{1/(2의 n제곱)}인가}

\(n\)큐비트 상태의 성분 수를 \(N=2^n\)이라고 하자.
연산자는 \(N\times N\) 행렬이고, 필요한 파울리 기저는 \(4^n=N^2\)개다.
각 큐비트마다 \(I,X,Y,Z\) 중 하나를 고르고 텐서곱한 행렬을
\textbf{파울리 문자열}이라고 부른다. 텐서곱은 각 큐비트에 지정한 작용을
함께 수행하는 다중 큐비트 행렬을 만드는 방법이다.

파울리 문자열 \(P,Q\)에는
\[
\operatorname{Tr}(P^\dagger Q)=
\begin{cases}N,&P=Q,\\0,&P\ne Q\end{cases}
\]
가 성립한다. 그러므로 전개와 계수는
\[
\boxed{\quad O=\sum_P a_PP,\qquad
a_P=\frac{\operatorname{Tr}(P^\dagger O)}{N},\quad N=2^n.\quad}
\]
파울리 문자열은 에르미트이므로 여기서 \(P^\dagger=P\)다.
한 큐비트에서는 2로, 두 큐비트에서는 4로 나눈다.
기저 행렬의 길이가 1이 아니라 \(\sqrt N\)이어서 이 분모가 생긴다.
길이까지 1로 맞춘 정규직교 기저를 원하면 \(P/\sqrt N\)를 쓴다.
이는 원래 파울리 행렬을 쓰는 위 계수식의 \(1/N\)과 구별해야 한다.

\section{파울리 문자열의 글자와 큐비트 번호}

두 큐비트에는 \(II,IX,IY,IZ,XI,\ldots,ZZ\) 등 16개 문자열이 있다.
이 노트와 Qiskit의 약속에서는 \textbf{가장 오른쪽 글자가 \(q_0\)에 작용한다}~\cite{ibm-order}.
예를 들어
\[
IZ=I_{q_1}\otimes Z_{q_0}=Z_0,\qquad
ZI=Z_{q_1}\otimes I_{q_0}=Z_1.
\]
주소 순서 \(00,01,10,11\)에서 대각 성분을 직접 적으면
\(IZ=\operatorname{diag}(1,-1,1,-1)\),
\(ZI=\operatorname{diag}(1,1,-1,-1)\)다.
오른쪽 비트가 1인 자리에 부호를 바꾸는지, 왼쪽 비트가 1인 자리에 바꾸는지 비교하면 된다.
\(Z_0\)처럼 아래첨자로 한 큐비트를 지정하면 나머지 큐비트에는 \(I\)가 있다는 뜻이다.
\(ZZ=Z_1Z_0\)에서는 서로 다른 두 큐비트에 작용하므로 곱으로도 적을 수 있다.

\begin{examplebox}{주소 하나로 순서를 확인하기}
\(\lvert q_1q_0\rangle=\lvert01\rangle\)에서는 \(q_0=1\), \(q_1=0\)이다.
따라서
\[
IZ\lvert01\rangle=-\lvert01\rangle,\qquad
ZI\lvert01\rangle=+\lvert01\rangle.
\]
\(IZ\)와 \(ZI\)를 바꾸면 다른 관측량이 된다.
회로 그림에서는 \(q_0\)가 윗줄에 놓일 수 있다. 문자열의 왼쪽을 회로의 윗줄로
외우기보다 큐비트 번호를 확인하는 편이 안전하다.
\end{examplebox}

세 큐비트가 일렬로 놓인 작은 모형도 적어 보자.
\[
\mathcal H=Z_0Z_1+Z_1Z_2-\tfrac12(X_0+X_1+X_2)
=IZZ+ZZI-\tfrac12(IIX+IXI+XII).
\]
\(ZZ\) 항은 이웃한 두 큐비트가 같은 부호인지에 따라 값을 주고,
\(X\) 항은 각 큐비트의 \(X\) 방향 성분을 읽는다.
이 \(\mathcal H\)를 에너지 관측량인 해밀토니안으로 해석할 수 있다.
여기서는 결합의 크기를 에너지 단위로 삼아 계수를 적었다.
아다마르 게이트 \(H\)와 혼동하지 않도록 이 해밀토니안은 \(\mathcal H\)로 표기한다.

\section{X와 Y를 읽으려면 측정 전에 무엇을 하나}

이제 행렬을 장치의 측정 절차와 연결하자.
표준적인 계산 기저 측정은 \(Z\)를 읽는다.
\(X\)나 \(Y\)를 읽을 때에는 측정 전에 상태를 회전시켜, 그 정보를 \(Z\) 측정에 맞춘다.
앞에서 \(H\)가 \(\lvert+\rangle,\lvert-\rangle\)을
\(\lvert0\rangle,\lvert1\rangle\)로 보낸 이유가 여기에 쓰인다.

측정 전에 적용하는 게이트를 \(V\)라고 하면 장치가 읽는 평균은
\[
\langle\psi\rvert V^\dagger ZV\lvert\psi\rangle
\]
다. 따라서 \(V^\dagger ZV=X\)가 되게 \(V=H\)를 고르면 \(X\)를 읽을 수 있다.
\(Y\)의 경우에는 \(S=\operatorname{diag}(1,i)\)를 사용한다.
\textbf{먼저 \(S^\dagger\), 그다음 \(H\)}를 적용한다.
\[
V=HS^\dagger,\qquad V^\dagger ZV=SHZHS^\dagger=Y.
\]
행렬곱은 오른쪽부터 상태에 작용한다. 그래서 실행 순서
\(S^\dagger\to H\)가 행렬로는 \(HS^\dagger\)다.

\lecturefigure{measurement-rotations}{기본 Z 측정으로 X, Y, Z를 읽는 세 절차. 회로의 시간은 왼쪽에서 오른쪽으로 흐른다. Y에서는 S의 켤레전치 게이트를 먼저 적용한다. 본문의 기저변환 항등식에서 구성한 자체 회로도다.}

\(\lvert+i\rangle\)에 \(S^\dagger\)를 적용하면 \(\lvert+\rangle\),
다시 \(H\)를 적용하면 \(\lvert0\rangle\)이 된다.
따라서 이 상태의 \(Y\) 측정값은 확실히 \(+1\)이다.
순서를 거꾸로 바꾸면 같은 측정이 되지 않는다.

여기서 세 종류의 기저가 만났다. 상태 기저는 상태벡터를 적는 기준이고,
파울리 기저는 행렬을 적는 기준이다. 측정 기저는 어떤 고유상태를 구별해 읽는가를 뜻한다.
\(X\)를 파울리 전개의 한 항으로 쓴다는 사실과,
\(X\)의 고유상태를 구별하도록 측정 장치를 준비한다는 사실은 이처럼 연결된다.

\section{여러 비트에서 하나의 파울리 측정값 얻기}

\(Z_1Z_0\)를 읽으려면 두 큐비트를 기본 측정한 뒤 두 부호를 곱하면 된다.
이 절에서는 측정 결과를 \(q_i\to c_i\)에 저장하고 \(b_1b_0\) 순서로 적는다.
그러면 한 번의 기록에서 얻는 값은
\[
p=(-1)^{b_1}(-1)^{b_0}=(-1)^{b_1+b_0}.
\]
00과 11은 \(+1\), 01과 10은 \(-1\)이다.
1의 개수가 짝수인지 홀수인지에 따라 부호가 정해져서 이를 패리티라고 부른다.

\lecturefigure{pauli-parity}{두 비트 기록을 ZZ의 측정값으로 바꾸는 표. 아래 빈도는 설명을 위해 만든 가상 기록이며 실험 데이터가 아니다. 부호별 횟수를 합쳐 평균을 구한다.}

예를 들어 같은 상태를 100번 새로 준비하고 측정한 기록이
00:37회, 11:43회, 01:12회, 10:8회였다고 가정하자.
이 가상 기록의 표본 평균은
\[
\widehat{\langle Z_1Z_0\rangle}
=\frac{(37+43)-(12+8)}{100}=0.6.
\]
0.6은 이 유한한 기록에서 얻은 추정값이다.
이상적인 독립 반복 측정만 가정하면 표준오차의 근삿값은
\(\sqrt{(1-0.6^2)/100}=0.08\)이다. 실제 장치에서는 여기에 측정과 게이트의 오류도 영향을 준다.

일반 문자열에서도 \(X\) 자리에는 \(H\), \(Y\) 자리에는
\(S^\dagger\) 다음 \(H\)를 적용한 뒤 해당 비트의 부호를 곱한다.
\(I\) 자리는 언제나 1이므로 그 비트를 곱에서 제외한다.
예를 들어 \(X_1Z_0\)는 \(q_1\)에만 \(H\)를 적용한 뒤 두 비트의 패리티로 읽는다.
클래식 비트에 다른 순서로 저장했다면 그 대응을 먼저 되짚어야 한다.

\begin{examplebox}{각각의 평균과 곱의 평균은 다르다}
벨 상태 \(\lvert\Phi^+\rangle=(\lvert00\rangle+\lvert11\rangle)/\sqrt2\)에서는
00과 11이 반반 나온다. 각 비트를 따로 보면
\(\langle Z_0\rangle=\langle Z_1\rangle=0\)이지만,
두 부호의 곱은 매번 \(+1\)이어서 \(\langle Z_1Z_0\rangle=1\)이다.
상관관계를 보려면 개별 평균을 곱하지 말고 각 기록의 곱을 평균해야 한다.
\end{examplebox}

\section{항의 평균을 더하는 일과 게이트를 잇는 일}

앞서 만든 관측량 \(O=2I+\tfrac12X-\tfrac14Y+\tfrac34Z\)는
각 파울리 항의 평균을 알아내면 전체 평균도 알 수 있다.
\[
\langle O\rangle=2+\tfrac12\langle X\rangle
-\tfrac14\langle Y\rangle+\tfrac34\langle Z\rangle.
\]
같은 상태를 준비하는 회로를 반복 사용하면서, 일부 반복에서는 \(X\), 다른 반복에서는
\(Y\), 나머지 반복에서는 \(Z\)를 읽으면 된다.
각 평균에 계수를 곱해 더한다. 이 절차가 파울리 전개를 측정에 이용하는 방법이다.

두 행렬이 가환한다는 말은 곱의 순서를 바꿔도 결과가 같다는 뜻이다.
\(X\)와 \(Z\)처럼 서로 가환하지 않는 관측량은 같은 상태의 측정값이
동시에 확실히 정해져 있지 않을 수 있다.
한 큐비트를 먼저 \(X\)로 측정한 뒤 \(Z\)로 측정하면 첫 측정이 상태를 바꾸므로,
두 번째 값을 원래 상태의 \(Z\) 측정 표본으로 취급할 수 없다.
위 절차에서는 \textbf{각 반복마다 원래 상태를 다시 준비한다}는 점이 중요하다.

전체 \(O\)의 고유기저로 직접 측정하는 방법과, 파울리 항을 여러 설정으로 나누어
평균을 조합하는 방법도 구별해야 한다. 후자는 기대값을 구하는 방법이다.
한 번의 기록에 가중치를 모두 더한 수가 반드시 \(O\)의 고윳값이 되지는 않는다.

측정 설정을 함께 쓸 수 있는 항은 묶을 수 있다.
예를 들어 \(ZI,IZ,ZZ\)는 같은 두 비트의 기본 측정 기록으로 모두 추정할 수 있다.
반면 \(XX\)와 \(ZZ\)는 전체 행렬로서는 가환하지만, 위와 같은 큐비트별 회전만 쓰는
단일 측정 설정은 서로 다르다. 이 둘을 함께 읽는 공동 고유기저 측정에는
추가적인 다중 큐비트 조작이 필요하다.

\begin{notebox}{합은 회로의 실행 순서가 아니다}
\(O=aX+bZ\)라는 식은 두 행렬을 가중합한다는 뜻이다.
\(X\) 게이트 다음 \(Z\) 게이트를 적용하면 행렬은 \(ZX\)가 된다.
\(aX+bZ\)와 \(ZX\)는 다른 연산이다. 파울리 전개식에 항이 보인다고
그 항을 그대로 차례로 게이트에 놓아서는 안 된다.
\end{notebox}

\section{네 무늬로 위상 표를 나누기}

처음의 네 칸 그림으로 돌아가자. 네 주소에 붙일 수를 다음 네 부호 무늬로 표현해 보자.
\[
\begin{aligned}
r_0&=(1,1,1,1),&r_1&=(1,-1,1,-1),\\
r_2&=(1,1,-1,-1),&r_3&=(1,-1,-1,1).
\end{aligned}
\]
서로 다른 두 행을 성분별로 곱해 더하면 0이 되고, 같은 행끼리 내적하면 4가 된다.
그러므로 네 행은 길이가 2인 직교 기저다.
\(f=(0,0,0,1)\)을 이 무늬로 나누면
\[
f=\frac14r_0-\frac14r_1-\frac14r_2+\frac14r_3.
\]
각 계수는 해당 무늬와 \(f\)의 내적을 4로 나눈 값이다.
예를 들어 \(r_1\cdot f=-1\)이므로 계수는 \(-1/4\)다.
네 항을 더하면 처음 세 주소에서는 0, 마지막 주소에서는 1이 된다.
이것이 여기서 사용하는 Walsh 전개다~\cite{welch}.

\lecturefigure{hadamard-patterns}{주소 순서 00, 01, 10, 11에서 본 네 Walsh 부호 무늬. 네 번째 주소의 표시값을 내적으로 읽어 계수 (1, -1, -1, 1)/4를 얻는다. 행의 길이와 계수는 본문의 계산을 사용했다.}

네 행을 모은 행렬을 \(R_4\)라고 하면 \(R_4R_4^\mathsf T=4I\)다.
양자 게이트로 쓰려면 벡터 길이를 보존해야 하므로 \(R_4/2=H\otimes H\)를 사용한다.
\(f\)의 고전 전개 계수에 붙는 \(1/4\)와, 양자 게이트의 정규화에 붙는 \(1/2\)는
서로 다른 역할이다. 일반적으로 주소 수가 \(N\)이면 각각 \(1/N\), \(1/\sqrt N\)가 된다.

\subsection{부호 무늬가 Z 문자열의 대각선이 된다}

같은 주소 순서에서 \(II,IZ,ZI,ZZ\)의 대각 성분은 각각
\(r_0,r_1,r_2,r_3\)다. 고전 무늬를 분해하던 계산이 곧 대각 행렬의 파울리 전개가 된다.
이번에는 부호를 바로 전개하는 대신, 부호를 만들어 낼 \textbf{위상 각도}를 전개하자.
검은 칸에 \(\pi\), 흰 칸에 0을 놓으면
\[
D=\pi\operatorname{diag}(0,0,0,1)
=\frac\pi4(II-IZ-ZI+ZZ).
\]
\(D\)는 실수 대각 에르미트 행렬이고 각도의 단위는 라디안이다.
위상 각도 \(\theta\)는 복소수 \(e^{i\theta}=\cos\theta+i\sin\theta\)로 진폭에 곱한다.
각도가 0이면 \(+1\), \(\pi\)이면 \(-1\)이 된다.
대각 행렬의 지수함수는 각 대각 성분에 지수함수를 적용한 행렬이므로,
이 \(D\)에서 원하는 게이트를 얻는다.
\[
U_f=e^{iD}=\operatorname{diag}(1,1,1,e^{i\pi})
=\operatorname{diag}(1,1,1,-1).
\]
\(D\)는 위상 각도를 담는 생성자이고, \(U_f\)는 상태에 실제로 작용하는 유니터리다.
해밀토니안으로 시간발전을 쓰는 관례 \(e^{-it\mathcal H/\hbar}\)와 부호가 달라 보일 수 있는데,
여기서는 처음에 정한 \(e^{+iD}\)를 끝까지 사용한다.

\section{위상 각도에서 정확한 게이트 회로로}

\(D\)의 네 항은 모두 대각 행렬이라 서로 가환한다.
즉 순서를 바꿔 곱해도 같다. 이 조건 덕분에 합의 지수함수를 정확히 나눌 수 있다.
\[
e^{iD}=e^{i\pi II/4}e^{-i\pi IZ/4}
e^{-i\pi ZI/4}e^{i\pi ZZ/4}.
\]
이 식에는 \(II\) 항도 있다. 그 항이 만드는 \(e^{i\pi/4}\)는
모든 진폭에 똑같이 붙는 공통 위상이다.
단독 상태의 측정 확률에는 영향을 주지 않지만, 전체 행렬이 정확히 같은 회로를
주장할 때에는 이 항까지 맞춰야 한다. 연산을 제어 게이트 안에 넣을 때에도
전체 위상이 상대위상으로 드러날 수 있다.

게이트 \(P(\theta)=\operatorname{diag}(1,e^{i\theta})\)를 쓰면 다음 회로가
공통 위상까지 포함해 \(U_f\)와 정확히 같다.
\begin{enumerate}
\item \(q_0,q_1\)에 각각 \(P(\pi/2)\)를 적용한다.
\item \(q_1\)을 제어, \(q_0\)를 표적으로 CX를 적용한다.
\item \(q_0\)에 \(P(-\pi/2)\)를 적용한다.
\item 같은 CX를 다시 적용해 주소를 원래대로 돌린다.
\end{enumerate}

\lecturefigure{phase-oracle}{네 칸 위상 오라클을 구현하는 P·CX 회로. 제어는 q1, 표적은 q0이며 시간은 왼쪽에서 오른쪽이다. 아래 식은 각 주소가 얻는 전체 위상을 보여준다. 위상 표와 게이트 행렬을 직접 계산해 구성했다.}

입력 비트를 \(x_1,x_0\)라고 하자. 첫 CX 뒤 표적 비트는
\(x_0\mathbin{\oplus}x_1\)가 된다. \(\oplus\)는 서로 다를 때 1인 배타적 논리합이다.
세 위상 게이트가 주는 총 각도는
\[
\frac\pi2x_0+\frac\pi2x_1-\frac\pi2(x_0\mathbin{\oplus}x_1)
=\pi x_0x_1.
\]
00,01,10에서는 0이고 11에서만 \(\pi\)다. 두 번째 CX는 비트를 복구하므로
네 입력 모두 주소는 그대로이고 원하는 부호만 남는다.

흔히 쓰는 \(R_z(\theta)=\operatorname{diag}(e^{-i\theta/2},e^{i\theta/2})\)와는
\(P(\theta)=e^{i\theta/2}R_z(\theta)\)라는 관계가 있다.
따라서 \(P\)를 \(R_z\)로 바꿀 때에는 공통 위상도 함께 추적해야 한다.

\begin{examplebox}{부호가 측정에서 드러나는 작은 후속 회로}
두 큐비트를 \(\lvert++\rangle\)로 준비하면 오라클 뒤 상태는
\((\lvert00\rangle+\lvert01\rangle+\lvert10\rangle-\lvert11\rangle)/2\)다.
바로 기본 측정하면 네 확률은 여전히 \(1/4\)이다.
여기에 \(q_0\)의 \(H\)만 더하면
\((\lvert00\rangle+\lvert11\rangle)/\sqrt2\)가 되어 00과 11만 반반 나온다.
상대 부호가 간섭을 거쳐 관측 가능한 상관관계로 바뀐 것이다.
\end{examplebox}

\subsection{X 항까지 있는 합에서는 조건을 다시 확인한다}

앞의 정확한 분리는 \(Z\) 문자열들이 서로 가환하기 때문에 가능했다.
앞서 본 이징 모형의 \(Z_1Z_0\)와 \(X_0\)는
\[
[Z_1Z_0,X_0]=Z_1Z_0X_0-X_0Z_1Z_0=2iZ_1Y_0\ne0
\]
이다. 따라서 일반적으로
\(e^{-it(A+B)}=e^{-itA}e^{-itB}\)라고 쓸 수 없다.
짧은 시간 간격으로 나누어 근사하는 방법도 있지만, 그때에는 시간 간격과 오차를
다루는 별도의 문제가 생긴다. 파울리 전개 자체는 항상 가능해도,
그 항들의 지수 게이트를 곧바로 이어 붙이는 절차가 항상 정확한 것은 아니다.

\section{Qiskit 코드로 적을 때의 짧은 연결}

이 절의 코드는 앞에서 손으로 정한 관측량을 Qiskit 자료형에 옮기는 예다.
공식 \texttt{SparsePauliOp} API를 대조했다~\cite{ibm-sparse}.
이번 개정에서는 SDK를 설치해 코드를 실행하지 않았다.
본문의 행렬·기대값은 별도의 작은 고전 행렬 계산으로 검산했으며,
아래에 SDK 실행 출력이나 QPU 결과를 제시하지 않는다.

\begin{verbatim}
from qiskit.quantum_info import SparsePauliOp

observable = SparsePauliOp.from_list([
    ("I", 2.0), ("X", 0.5),
    ("Y", -0.25), ("Z", 0.75),
])

ising3 = SparsePauliOp.from_sparse_list([
    ("ZZ", [0, 1], 1.0),
    ("ZZ", [1, 2], 1.0),
    ("X", [0], -0.5),
    ("X", [1], -0.5),
    ("X", [2], -0.5),
], num_qubits=3)
\end{verbatim}

\texttt{from\_list}는 전체 문자열과 계수의 목록을 받고,
\texttt{from\_sparse\_list}는 작용할 큐비트 번호를 따로 받는다.
이 예제의 \texttt{ZZ}는 같은 글자 두 개이므로 번호별 작용이 분명하다.
서로 다른 글자를 번호 목록과 함께 쓸 때에도 각 글자가 어느 번호에 대응하는지
API의 약속을 확인해야 한다.

행렬을 이미 알고 있다면 다음처럼 파울리 전개를 요청할 수도 있다.
\begin{verbatim}
import numpy as np

matrix = np.array([
    [2.75, 0.5 + 0.25j],
    [0.5 - 0.25j, 1.25],
], dtype=complex)
decomposed = SparsePauliOp.from_operator(matrix)
\end{verbatim}
수학적으로 앞의 \(2I+\tfrac12X-\tfrac14Y+\tfrac34Z\)와 같은 연산자다.
다만 코드가 항을 어떤 순서로 출력하는지는 관측량의 의미와 무관하다.

\texttt{SparsePauliOp}에서 ``sparse''는 필요한 파울리 항만 저장한다는 뜻이다.
모든 입력이 적은 항으로 표현된다는 보장은 아니다.
\(n\)큐비트의 조밀한 행렬은 \(4^n\)개 성분을 가지며,
일반적인 파울리 전개도 최대 \(4^n\)항이 필요하다.
작은 행렬을 변환하는 이 예제를 큰 임의 행렬의 효율적인 처리법으로 확대해서는 안 된다.

\subsection{Estimator에 주는 두 가지 입력}

기대값을 구하려면 준비할 상태와 읽을 관측량이 함께 필요하다.
회로는 상태 준비를, \texttt{SparsePauliOp}는 관측량을 기술한다.
Estimator 계열 인터페이스는 이 둘을 받아 기대값을 계산하거나 추정하는 역할을 한다.
구현에 따라 필요한 변환·측정과 실행 방식이 달라지므로,
측정 전 회전을 직접 넣는 예와 Estimator 사용을 혼동해 중복 적용하지 않는다.

예를 들어 \texttt{StatevectorEstimator}는 이상적인 유니터리 회로의 상태벡터로
정확한 기대값을 계산할 수 있다~\cite{ibm-estimator}.
실물 장치에서 얻는 유한한 샷의 평균과는 실행 조건이 다르다.
관측량 객체를 만들었다는 사실만으로 회로를 실행하거나 실험값을 얻은 것은 아니다.

\section{같은 도구가 어디까지 이어지는가}

네 칸의 위상 각도를 Walsh 무늬로 나누고, 그 무늬를 \(Z\) 문자열로 바꾸면
대각 위상 회로를 유도할 수 있었다.
일반적인 \(n\)큐비트 대각 행렬도 주소 \(x\)마다 위상 각도 \(\theta_x\)를 정해
\[
D=\operatorname{diag}(\theta_x)=\sum_{s\in\{0,1\}^n}d_sZ^s,
\qquad d_s=\frac1{2^n}\sum_x\theta_x(-1)^{s\cdot x}
\]
로 쓸 수 있다. \(Z^s\)는 \(s_j=1\)인 큐비트에만 \(Z\)를 놓은 문자열이다.
\(s\cdot x\)의 홀짝이 그 주소에서의 부호를 정한다.
일반 파울리 계수식이 대각 행렬에서는 고전 Walsh 계수식으로 줄어든 것이다.

이 계수 계산은 고전 컴퓨터에서 수행할 수 있다.
위상 표 자체가 \(2^n\)개 수를 담으면 그 입력과 전개가 이미 클 수 있다.
회로가 짧아지는지는 살아남는 계수의 수와 구조, 구현할 연결과 게이트에 달려 있다.
기저로 표현할 수 있다는 사실만으로 양자 이득이나 짧은 회로가 보장되지는 않는다.

마지막으로, 여기서 쓴 \(H^{\otimes n}\)은 차수가 \(2^n\)인 Sylvester형
아다마르 행렬을 정규화한 것이다. 조합론에서는 성분이 \(\pm1\)이고 서로 직교하는
행을 가진 더 넓은 종류의 아다마르 행렬을 연구한다~\cite{egan}.
4의 배수인 모든 차수에 그런 행렬이 존재하느냐는 일반 문제와,
여기서 명시적으로 만들 수 있는 \(2^n\)차 행렬의 존재는 구별해야 한다.
이 노트의 회로 구성은 그 일반 추측의 해결을 요구하지 않는다.

\section{손으로 확인하는 다섯 문제}

\begin{questionbox}{1. 같은 기본 확률, 다른 X 평균}
\(\lvert+\rangle\)와 \(\lvert-\rangle\)을 기본 측정하면 확률이 같다.
각 상태의 \(\langle X\rangle\)는 얼마인가? 어떻게 읽을 수 있는가?
\end{questionbox}
\begin{answerbox}{풀이}
각각 \(+1,-1\)이다. \(H\)를 적용하면 각각 \(\lvert0\rangle,\lvert1\rangle\)이 되어,
기본 측정 뒤 0을 \(+1\), 1을 \(-1\)로 읽으면 된다.
\end{answerbox}

\begin{questionbox}{2. 행렬을 네 파울리 항으로 나누기}
\(A=\begin{pmatrix}3&1-i\\1+i&1\end{pmatrix}\)를
\(I,X,Y,Z\)로 표현하고 \(\lvert+i\rangle\)에서의 기대값을 구하라.
\end{questionbox}
\begin{answerbox}{풀이}
대각 성분의 평균과 차이에서 \(a=2,d=1\), 비대각 성분에서 \(b=1,c=1\)이다.
따라서 \(A=2I+X+Y+Z\)이고 기대값은 \(2+0+1+0=3\)이다.
\end{answerbox}

\begin{questionbox}{3. 문자열 순서와 비트의 부호}
\(\lvert q_1q_0\rangle=\lvert10\rangle\)에서
\(IZ,ZI,ZZ\)의 기대값을 구하라.
\end{questionbox}
\begin{answerbox}{풀이}
\(q_0=0,q_1=1\)이므로 각각 \(+1,-1,-1\)이다.
\(ZZ\)는 두 부호의 곱이다.
\end{answerbox}

\begin{questionbox}{4. 가상 기록에서 XX 평균 얻기}
각 큐비트에 \(H\)를 적용한 뒤 기본 측정했다.
같은 번호의 클래식 비트에 저장한 기록이 00:30회, 11:40회, 01:20회, 10:10회다.
원래 상태의 \(XX\) 평균 추정값은 얼마인가?
\end{questionbox}
\begin{answerbox}{풀이}
\(H\) 덕분에 각 비트가 \(X\) 측정에 대응한다.
짝수 패리티 70회, 홀수 패리티 30회이므로 \((70-30)/100=0.4\)다.
\end{answerbox}

\begin{questionbox}{5. 마지막 칸만 바꾸는 위상}
\(D_\varphi=\varphi\operatorname{diag}(0,0,0,1)\)의 파울리 전개를 쓰고,
앞의 P·CX 회로에서 어떤 각도를 바꾸면 되는지 설명하라.
\end{questionbox}
\begin{answerbox}{풀이}
\(D_\varphi=(\varphi/4)(II-IZ-ZI+ZZ)\)다.
두 첫 P 게이트의 각도를 \(\varphi/2\), 가운데 P 게이트의 각도를
\(-\varphi/2\)로 바꾸면 총 위상은 \(\varphi x_0x_1\)가 된다.
따라서 \(e^{iD_\varphi}\)를 정확히 구현한다.
\end{answerbox}

\begin{thebibliography}{7}
\bibitem{ibm-observables} IBM Quantum, \href{https://quantum.cloud.ibm.com/docs/en/qiskit/specify-observables-pauli}{Specify observables in the Pauli basis}, IBM Quantum Documentation (확인 2026-11-28; 문서의 코드 환경 \texttt{qiskit[all]\textasciitilde{}=2.6.1}).
\bibitem{preskill} J. Preskill, \href{https://www.preskill.caltech.edu/ph219/chap2_15.pdf}{Lecture Notes for Ph219/CS219: Quantum Information, Chapter 2: Foundations I: States and Ensembles} (July 2015).
\bibitem{ibm-order} IBM Quantum, \href{https://quantum.cloud.ibm.com/docs/en/qiskit/bit-ordering}{Bit-ordering in the Qiskit SDK}, IBM Quantum Documentation (확인 2026-11-28).
\bibitem{welch} J. Welch, D. Greenbaum, S. Mostame, and A. Aspuru-Guzik, \href{https://doi.org/10.1088/1367-2630/16/3/033040}{Efficient quantum circuits for diagonal unitaries without ancillas}, New J. Phys. \textbf{16}, 033040 (2014).
\bibitem{ibm-sparse} IBM Quantum, \href{https://quantum.cloud.ibm.com/docs/en/api/qiskit/qiskit.quantum_info.SparsePauliOp}{SparsePauliOp}, Qiskit API Reference (확인 2026-11-28).
\bibitem{ibm-estimator} IBM Quantum, \href{https://quantum.cloud.ibm.com/docs/en/api/qiskit/qiskit.primitives.StatevectorEstimator}{StatevectorEstimator}, Qiskit API Reference (확인 2026-11-28).
\bibitem{egan} R. Egan, I. S. Kotsireas, S. London, and P. Ó Catháin, \href{https://arxiv.org/abs/2610.27175v2}{Recent progress in the construction of Hadamard matrices}, arXiv:2610.27175v2 (2026).
\end{thebibliography}

\section*{작성·도해·검산과 판본 안내}
\addcontentsline{toc}{section}{작성·도해·검산과 판본 안내}

책임 편집자는 김현중이다. 이 v3는 사용자가 지정한 학습 주제와 개정 방향에 따라
편집부(AEDAUT)가 OpenAI Codex의 집필·감수 보조로 전체 설명을 새로 정리한 한국어 학습 노트다.
기존 v2의 네 칸 위상 오라클, Walsh 무늬, 기저의 구별과 P·CX 회로를 유지하면서,
관측량의 뜻, 파울리 계수의 유도, 측정 회전과 패리티, 기대값 계산을 원래의 작은 예제에 연결했다.
IBM 문서의 번역본이 아니며 참고문헌과 독립 계산을 바탕으로 구성했다.
여섯 설명도는 본문의 행렬·규칙으로 새로 그린 벡터 도해다. 외부 그림을 복제하지 않았다.

기존 v2 공개판과 그 판본의 필기 연결은 보존한다. 같은 렉처 주소에서 판본에 맞는 공책을 사용할 수 있다.
수신 원판의 집필 모델과 원저자의 별도 검산 기록은 확인되지 않았다.
이전 공개 편집의 기록과 이번 v3의 집필·검토 기록은 구분한다.
사람 독자의 이해도 시험이나 사람의 문장별 감수를 수행한 것으로 표시하지 않는다.
이번 작업에서 확인한 에이전트 식별명과 역할은 \texttt{root}(집필·편집 보조),
\texttt{ftqc\_science}(원전 대조·과학 감수), \texttt{ftqc\_pdf}(도해·PDF 제작),
\texttt{pqc\_archive\_structure}(독자 흐름·판본 연속성·발행 검수)다.

이번 v3에서는 별도 작성한 Python 표준 라이브러리 계산으로 이름을 붙인 27개 수학 항목을 검산했다.
세 상태의 기대값, 파울리 계수와 재구성, 측정 회전, 비트 순서와 패리티,
벨 상태의 상관관계, Walsh 계수, 공통 위상을 포함한 P·CX 회로,
비가환 지수 분해의 반례, 다섯 문제의 풀이와 세 큐비트 모형을 확인했다.
최대 행렬 크기는 \(8\times8\), 항등식 검산의 최대 잔차는 약 \(8.9\times10^{-16}\),
허용 기준은 \(10^{-12}\)였다. 비가환 반례에서는 두 식의 차이가 0이 아님을 따로 확인했다.
27은 검산 묶음의 수이며 실험 횟수가 아니다.
이전 판본의 검산을 재사용한 것으로 세지 않고 이번에 직접 다시 계산했다.
SDK·QPU·실물 실험과 원저자의 별도 검산 프로그램은 실행하지 않았다.
가상 측정 기록은 본문에 명시한 교육용 예다.


\AIModelDisclosure

